% ============================================================
%  BDS-SMC2 Node — Comprehensive Report (LaTeX rendition)
%  Compile:  pdflatex BDS_SMC2_Node_Comprehensive_Report.tex
%            (run twice for the table of contents / links)
%  Figures expected relative to this file: figures/  and  _photos/
% ============================================================
\documentclass[11pt,a4paper]{article}

\usepackage[T1]{fontenc}
\usepackage[utf8]{inputenc}
\usepackage{lmodern}
\usepackage[margin=1in]{geometry}
\usepackage{microtype}
\usepackage{graphicx}
\usepackage{float}
\usepackage{enumitem}
\usepackage{booktabs}
\usepackage{array}
\usepackage{tabularx}
\usepackage{longtable}
\usepackage{seqsplit}
\usepackage{xcolor}
\usepackage{tcolorbox}
\usepackage{pifont}
\usepackage{fancyvrb}
\usepackage{caption}
\usepackage{hyperref}
\hypersetup{colorlinks=true,linkcolor=band,urlcolor=blue,citecolor=band}

\definecolor{band}{HTML}{0D2B45}
\definecolor{accent}{HTML}{D9822B}
\definecolor{good}{HTML}{2E8B57}
\definecolor{warn}{HTML}{B03A4B}
\definecolor{rowbg}{HTML}{F2F5F8}

% symbol shortcuts
\newcommand{\yes}{\textcolor{good}{\ding{51}}}
\newcommand{\no}{\textcolor{warn}{\ding{55}}}
\newcommand{\partialmark}{\textcolor{accent}{$\circ$}}

% callout box
\newtcolorbox{callout}[1][]{colback=band!4,colframe=band,boxrule=0.6pt,
  arc=2pt,left=6pt,right=6pt,top=4pt,bottom=4pt,#1}

% verbatim style for the ASCII diagrams
\DefineVerbatimEnvironment{ascii}{Verbatim}%
  {fontsize=\scriptsize,frame=single,framesep=4pt,rulecolor=\color{band!40}}

\captionsetup{font=small,labelfont=bf}
\setlength{\parskip}{0.5em}
\setlength{\parindent}{0pt}
\renewcommand{\arraystretch}{1.25}

\title{\vspace{-1.2em}\textbf{BDS--SMC2 Node --- Comprehensive Report \textit{(Upgraded)}}}
\author{Letsoalo Maile \, $\cdot$ \, WP5 (BDS Communication) \, $\cdot$ \, BDS--SMC2 Node}
\date{\today}

\begin{document}
\maketitle
\thispagestyle{empty}

\begin{callout}
\textbf{Scope of this document:} a single, self-contained account of the node --- what it
addressed, the sequence in which it was built, what it could do \emph{before} versus
\emph{now}, justification of every figure and statistic, every contradiction found and how
it resolves, the completed group integration, and a precise statement of work done
\emph{within} and \emph{beyond} my assigned task --- and the path to publication.

\smallskip
\textit{Every number in this report was regenerated from the committed scripts in
\texttt{python/} against the CSVs in \texttt{data/}. Nothing is asserted without a source.
Reproduction commands are in Appendix~A.}
\end{callout}

\tableofcontents
\vspace{1em}
\hrule

% ============================================================
\section*{0. Executive Summary}
\addcontentsline{toc}{section}{0. Executive Summary}

The BDS-SMC2 node began as a 64-bit, two-field, $\sim$11-metre coordinate transmitter and
is now a 112-bit, six-field, $\sim$1-centimetre \textbf{software-complete rescue pipeline}
that is \textbf{integrated and verified inside the group's ROS\,2 system}. The headline
compression result ($\sim$58\,\% smaller than ASCII) held through a full payload redesign,
which is evidence of a robust finding rather than a tuned one. The work assigned to me
(WP5: a simulation emulator + parser + injector) is \textbf{complete and merged}; on top of
it I delivered an unrequested hardware bring-up, a 232-transmission field-reliability
campaign, a latency baseline, and original encoding research. The one claim still unproven
--- physical-layer transmission of the 112-bit frame --- is honestly scoped as the
validation boundary, not hidden. \textbf{This is a successful WP5 delivery with a
substantial research contribution attached.}

% ============================================================
\section{What the Node Was Addressing}

The five-module team project (Beihang, ``UAV Emergency Rescue with BeiDou Short Message
Communication'') targets a disaster scenario where terrestrial networks collapse and
BeiDou-3 short messaging is the only surviving channel. \textbf{WP5 --- my module --- owns
the communications backbone:} take a survivor's RTK-corrected coordinate, carry it through
BeiDou short messaging, decode it at the ground station, and inject it into the autonomous
mission so the UAV can act on it.

The node specifically addresses four questions:
\begin{enumerate}[leftmargin=1.6em,itemsep=1pt]
  \item \textbf{Capacity} --- can a complete rescue payload fit one short message?
  \item \textbf{Reliability} --- does delivery survive real propagation environments?
  \item \textbf{Latency} --- does it arrive fast enough to coordinate a rescue?
  \item \textbf{Integration} --- can the decoded message drive an autonomous mission rather
        than a human-read display?
\end{enumerate}

% ============================================================
\section{Sequence of Events (chronological, git-backed)}

{\small
\begin{longtable}{@{}p{2.0cm} p{2.5cm} p{6.4cm} p{2.2cm}@{}}
\toprule
\textbf{Date} & \textbf{Commit} & \textbf{Event} & \textbf{Stage} \\
\midrule
\endhead
2026-06-06 & \texttt{83435ce} & Initial node: 64-bit payload, 2 fields, 4 dp & \textbf{v1 born} \\
2026-06-06 & \texttt{33ee9af} & Removed Gap 5 (AES); added analysis scripts + field directives & scope trim \\
2026-06-08 & \texttt{da015a8} & \textbf{Hardware day 1} --- T3 firmware fix + first real field data & hardware live \\
06-08$\to$09 & \texttt{1a55e9d} $\to$ \texttt{bd8e2d4} & Gap 3 field campaign: indoor, open-sky, canopy, urban-canyon --- \textbf{232 valid TX, 12 locations} & empirical data \\
2026-06-13 & \texttt{dcdadc7} & \textbf{64-bit $\to$ 112-bit} rescue payload (6 fields, 7 dp) + receive chain & \textbf{v2 born} \\
2026-06-13 & \texttt{dd18d37}, \texttt{b7c5b5b} & Paper 1 sections I--VIII drafted & writing \\
2026-06-13 & \texttt{137ef76} & \texttt{portal\_reader} live connection to BeiDou portal & receive live \\
2026-06-14 & \texttt{a752909}, \texttt{95dbf7e} & Auto-sender CSV-watch pipeline + retry policy & automation \\
2026-06-16 & \texttt{a0817d2} & Sim + GCS + virtual BeiDou link + dashboard (Objective 5 demo) & demo \\
group repo & \texttt{e37098e} $\to$ \texttt{1065eaa} & \textbf{Integrate all 5 modules (Stage-1)} $\to$ beidou node + 112-bit decode + \textbf{integration verify kit} & \textbf{integ. merged} \\
\bottomrule
\end{longtable}
}

\textbf{Reading of the sequence:} the field-reliability campaign (8--9 June)
\emph{predates} the 112-bit upgrade (13 June). That ordering matters and is referenced
repeatedly below --- it is the reason the empirical data uses the ASCII payload.

% ============================================================
\section{The Node --- Before vs Now (Capability Evolution)}

\subsection{What it COULD do (v1, 06-06)}
\begin{itemize}[leftmargin=1.4em,itemsep=1pt]
  \item Pack two coordinates (lat, lon) into \textbf{64 bits} at \textbf{4 dp ($\sim$11 m)}.
  \item Demonstrate \textbf{$\sim$57.9\,\%} size reduction versus an ASCII baseline (184 bit).
  \item Nothing else: no altitude, no uncertainty, no triage, no survivor ID; no receive
        chain; no field data; no mission integration.
\end{itemize}

\subsection{What it CAN do now (v2, current)}
\begin{itemize}[leftmargin=1.4em,itemsep=1pt]
  \item Pack \textbf{six} operationally complete fields (lat, lon, \textbf{alt, uncertainty
        R, priority, survivor ID}) into \textbf{112 bits} at \textbf{7 dp ($\sim$1 cm)},
        decoding \textbf{bit-for-bit} against ground truth.
  \item Decode three wire formats behind one interface: ASCII, 112-bit binary, legacy 64-bit.
  \item Carry a \textbf{232-transmission field-reliability} result and a \textbf{30-sample
        latency} baseline.
  \item Pull messages down through a \textbf{live portal reader}, and \textbf{publish the
        rescue trigger into the group ROS\,2 mission system} (Section~6).
\end{itemize}

\subsection{Side-by-side (every row reproduced)}
\begin{longtable}{@{}p{3.1cm} p{2.6cm} p{3.4cm} p{3.6cm}@{}}
\toprule
\textbf{Capability} & \textbf{v1 (before)} & \textbf{v2 (now)} & \textbf{Source} \\
\midrule
\endhead
Payload size & 64 bit & \textbf{112 bit} & \texttt{decode\_binary.py} \texttt{">ii"}$\to$\texttt{">iihHBB"} \\
Fields & 2 & \textbf{6} & Table 1, report \\
Precision & 4 dp ($\sim$11 m) & \textbf{7 dp ($\sim$1 cm)} & \texttt{*10000}$\to$\texttt{*10000000} \\
gap1 ASCII baseline & 184 bit & 264 bit & \texttt{gap1\_compression.csv} \\
gap1 reduction & $\sim$57.9\,\% & \textbf{57.6\,\%} & both eras --- \emph{stable} \\
gap6 ASCII/Bin/Huff & 384/128/192 & \textbf{368/128/184} & \texttt{gap6\_telemetry.csv} \\
Field reliability & none & \textbf{232 TX, $\ge$93.7\,\% Wilson LB} & \texttt{gap3\_analysis.py} \\
Latency & none & \textbf{2574.5 ms, n=30} & \texttt{gap2\_analysis.py} \\
Receive + mission & none & \textbf{portal $\to$ decode $\to$ ROS 2 trigger} & \texttt{gcs/}, group node \\
\bottomrule
\end{longtable}

\textbf{Defensible one-line claim:} \emph{the node moved from a 2-field, 11-metre,
transmit-only prototype to a 6-field, centimetre-precision, integrated rescue pipeline ---
while the core $\sim$58\,\% efficiency result survived the redesign.}

\subsection{Clear Illustration --- The Role I Was Tasked To Do}
\begin{ascii}
                THE 5-MODULE UAV RESCUE SYSTEM (Beihang team)
   +--------+--------+--------+--------+-----------------------------+
   |  WP1   |  WP2   |  WP3   |  WP4   |      WP5  =  MY TASK        |
   |  RTK   |  QGC   |  YOLO  |  Path  |   BDS COMMUNICATION         |
   | Tevin  | Yvonne |  Afiq  |  Aman  |      Letsoalo Maile         |
   +--------+--------+---+----+--------+--------------+--------------+
                         |                            |
                         |  consume                   | MY CONTRACT (WP5):
                         v                            v
        /target/emergency_coordinate  <-------  survivor coord
                  (ROS 2 topic)                     |
                                                    v
        +----------------------- MY WP5 PIPELINE ------------------------+
        |  encode  ->  BDS short message  ->  decode  ->  publish        |
        |  (Class+PNT)     (emulated/real)     (parse)    (inject PNT)   |
        +---------------------------------------------------------------+
   Required of me: a SIMULATION emulator + parser + injector, via ROS 2.
\end{ascii}

\subsection{Clear Illustration --- The Node BEFORE vs NOW}
\begin{ascii}
  +----------------------- BEFORE (v1, 06-06) -----------------------+
  |  [lat, lon]  ->  64-bit pack (4 dp ~ 11 m)  ->  TRANSMIT  ->  X   |
  |  - 2 fields, no altitude / uncertainty / triage / ID             |
  |  - no receive chain   - no field data   - no mission integration |
  |  RESULT: a transmitter that could not carry a complete rescue,   |
  |          and had nowhere to deliver it.                          |
  +------------------------------------------------------------------+
                                 ||  64->112-bit upgrade + field campaign
                                 vv
  +------------------------- NOW (v2, current) ----------------------+
  |  [lat,lon,alt,R,priority,ID] -> 112-bit pack (7 dp ~ 1 cm)       |
  |   BDS short message -> PORTAL -> portal_reader -> DECODE ->      |
  |       -> publish /target/emergency_coordinate -> UAV MISSION OK  |
  |  - 6 complete fields, centimetre precision, bit-exact round-trip |
  |  - 232-TX field reliability (>=93.7% Wilson)  - 2.57 s latency   |
  |  - INTEGRATED + verified in the group ROS 2 system               |
  |  RESULT: a software-complete, integrated, empirically-backed     |
  |          rescue pipeline.                                        |
  +------------------------------------------------------------------+
\end{ascii}

% ============================================================
\section{Figures --- Justified Against Data}

All six figures the report embeds exist, regenerate, and match their captions.

\begin{callout}
\textbf{Baseline-disclosure note (resolves C4):} the ASCII comparison baseline was
deliberately re-specified mid-project from the early 2-field/4-dp string (184 bit) and
7-field/4-dp telemetry (384 bit) to the \textbf{operational \texttt{\$CCTXM} rescue format}
(264 bit / 368 bit). All current figures use the operational baseline; where an older
number appears it is the superseded baseline, flagged as such --- the change is a
re-specification, not a silent re-tuning.
\end{callout}

{\small
\begin{longtable}{@{}c >{\raggedright\arraybackslash}p{3.7cm} p{4.4cm} >{\raggedright\arraybackslash}p{3.4cm}@{}}
\toprule
\textbf{Fig} & \textbf{File} & \textbf{Claim} & \textbf{Verified against} \\
\midrule
\endhead
1 & \texttt{\seqsplit{bds\_node\_workflow.png}} & End-to-end pipeline (encode $\to$ satellite $\to$ portal $\to$ decode $\to$ ROS 2) & matches \texttt{gcs/} + architecture \\
2 & \texttt{\seqsplit{gap1\_encoding\_comparison.png}} & 264 bit $\to$ 112 bit, \textbf{$-$57.6\,\%} & \texttt{\seqsplit{gap1\_compression.csv}} \\
3 & \texttt{\seqsplit{gap6\_telemetry\_comparison.png}} & ASCII 368 / Huffman 184 / Binary 112; 210-bit limit line & \texttt{\seqsplit{gap6\_telemetry.csv}} (see C3) \\
4 & \texttt{\seqsplit{gap3\_success\_rate.png}} & Per-environment delivery, Wilson 95\,\% intervals & \texttt{\seqsplit{gap3\_analysis.py}} \\
5 & \texttt{\seqsplit{gap3\_location\_breakdown.png}} & All 12 locations 100\,\% in-sample & \texttt{\seqsplit{gap3\_analysis.py}} \\
6 & \texttt{\seqsplit{fig\_gap2\_cdf.png}} & Latency CDF, all $<$ 5 s & \texttt{\seqsplit{gap2\_latency\_ascii\_baseline.csv}} \\
\bottomrule
\end{longtable}
}

\textbf{Key statistics reproduced live in this session:}
\begin{itemize}[leftmargin=1.4em,itemsep=1pt]
  \item Encoding: 368 / 184 / 112 bit; 264-bit rescue ASCII $\to$ 57.6\,\%;
        69.6\,\% vs full telemetry. \yes
  \item Reliability: 232 valid (61 / 57 / 57 / 57); $\chi^2$ = 0.000, df = 3, p = 1.0000;
        Fisher pairwise all 1.0000; Wilson LB 93.7\,\% per-env, 98.4\,\% pooled. \yes
  \item Latency: mean \textbf{2574 ms}, SD \textbf{1095 ms}, n = 30, max $<$ 4.5 s. \yes
  \item Round-trip: T001--T006 decode bit-exact at 7 dp. \yes
\end{itemize}

\begin{figure}[H]\centering
  \includegraphics[width=0.85\textwidth]{figures/bds_node_workflow.png}
  \caption{Fig.~1 --- End-to-end pipeline: encode $\to$ satellite $\to$ portal $\to$ decode $\to$ ROS~2.}
\end{figure}
\begin{figure}[H]\centering
  \includegraphics[width=0.49\textwidth]{figures/gap1_encoding_comparison.png}\hfill
  \includegraphics[width=0.49\textwidth]{figures/gap6_telemetry_comparison.png}
  \caption{Fig.~2 (left) encoding comparison 264$\to$112 bit; Fig.~3 (right) telemetry comparison.}
\end{figure}
\begin{figure}[H]\centering
  \includegraphics[width=0.49\textwidth]{figures/gap3_success_rate.png}\hfill
  \includegraphics[width=0.49\textwidth]{figures/gap3_location_breakdown.png}
  \caption{Fig.~4 (left) per-environment delivery with Wilson intervals; Fig.~5 (right) 12-location breakdown.}
\end{figure}
\begin{figure}[H]\centering
  \includegraphics[width=0.7\textwidth]{figures/fig_gap2_cdf.png}
  \caption{Fig.~6 --- Latency CDF; all transmissions under 5 s.}
\end{figure}

% ============================================================
\section{Contradictions Found --- and How Each Resolves}

Found during a line-by-line audit. Ranked. Each is either a free fix or an honest-framing item.

\subsection*{Tier C --- documentation/consistency (cheap, fix before submission)}
\begin{itemize}[leftmargin=1.4em,itemsep=2pt]
  \item \textbf{C1 --- README said ``Gap 3 = Pending hardware.''} Stale: git shows Gap 3
        collected 06-09. \emph{Resolution: corrected to ``Complete (ASCII payload, 232 TX).''}
  \item \textbf{C2 --- ``strictly more information than Huffman'' is false.} The 112-bit
        binary drops battery/mode/flags/timestamp that the 368-bit ASCII/Huffman carry.
        \emph{Resolution: reword to ``comparable size for a different, rescue-optimised field
        set''; lead with the clean 264$\to$112 comparison.}
  \item \textbf{C3 --- Figure 3 draws a hard ``210-bit limit''} the prose calls
        ``unmeasured.'' \emph{Resolution: relabel the figure line ``indicative'' or commit
        the prose --- not both.}
  \item \textbf{C4 --- the comparison baseline moved} (184$\to$264 bit, 384$\to$368 bit)
        without being flagged. \emph{Resolution: one sentence noting the re-specification to
        the operational \texttt{\$CCTXM} format.}
\end{itemize}

\subsection*{Tier B --- bounds on the existing results (honest framing)}
\begin{itemize}[leftmargin=1.4em,itemsep=2pt]
  \item \textbf{B1 --- zero failures everywhere $\to$ cannot resolve true reliability above
        $\sim$94\,\%} ($\chi^2$ test degenerate by construction). \emph{Handled by the
        Wilson-lower-bound framing; keep it.}
  \item \textbf{B2 --- ``twelve locations'' share one hardcoded coordinate}
        (30.4196, 120.2977); no genuine fade diversity. \emph{Disclose as emulated input;
        propose deep-fade sites as future work.}
\end{itemize}

\subsection*{Tier A --- the validation boundary (needs the lab)}
\begin{itemize}[leftmargin=1.4em,itemsep=2pt]
  \item \textbf{A1 --- the 112-bit frame has never been transmitted on hardware.} Software
        round-trip + sim only. \emph{Honestly scoped; one hardware TX closes it.}
  \item \textbf{A2 --- the reliability + latency campaigns used the ASCII payload.}
        \emph{A chronological artifact (field test predates the upgrade), which makes the
        reliability result payload-independent --- the more general claim.}
\end{itemize}

\textbf{Net:} no fabrication anywhere; the contradictions are documentation lag (C),
statistical honesty (B), and a clearly-marked hardware boundary (A).

% ============================================================
\section{Integration Status --- COMPLETE (group repo evidence)}

The BDS node is \textbf{merged into the group ROS\,2 system} and verified.

\begin{longtable}{@{}p{3.4cm} p{10.4cm}@{}}
\toprule
\textbf{Evidence} & \textbf{Detail} \\
\midrule
\endhead
Node & \texttt{beidou\_publisher\_node.py} imports \texttt{interfaces.msg/EmergencyCoordinate}, creates a \textbf{latched publisher} on \texttt{/target/emergency\_coordinate} \\
Shared contract & \texttt{interfaces/msg/EmergencyCoordinate.msg} (header, lat, lon, source\_id, raw\_message) \\
Consumers wired & \texttt{qgc\_control} and \texttt{path\_planning} subscribe to \texttt{/target/emergency\_coordinate} \\
Verify kit & \texttt{verify\_integration.sh} runs 3 checks; expected ``ALL CHECKS PASSED'' \\
Merge history & \texttt{e37098e} (Integrate all 5 modules) $\to$ \texttt{1065eaa} (beidou node + 112-bit decode + verify kit) \\
Shared datum & \texttt{af76d5c} (single Zurich datum across all modules) --- contract-compliant \\
\bottomrule
\end{longtable}

\textbf{The decoded coordinate flows end-to-end into mission planning.} The four extra
rescue fields (alt, R, priority, survivor ID) are decoded and \textbf{logged as a documented
enhancement talking-point}, not yet in the \texttt{.msg} --- an optional future extension,
\textbf{not} an unfinished integration.

% ============================================================
\section{My Contribution --- Within Scope vs Beyond Scope}

\subsection{Within my assigned scope (WP5, from the opening presentation)}
Defined on deck slides 4/8/9: \emph{BDS-SMC emulator node $\cdot$ coordinate message parsing
$\cdot$ QGC interface validation $\cdot$ auto-inject parsed BDS messages $\to$ precise PNT
$\cdot$ integrate via ROS in Phase 3.}

\begin{longtable}{@{}p{5.2cm} p{6.4cm} p{2.0cm}@{}}
\toprule
\textbf{WP5 deliverable} & \textbf{Delivered} & \textbf{Status} \\
\midrule
\endhead
BDS-SMC emulator node & \texttt{sim/bds\_module\_emulator.py}, \texttt{sim/virtual\_portal.py} & \yes \\
Coordinate message parsing & \texttt{decode\_ascii.py} / \texttt{decode\_binary.py} / \texttt{gcs/decoder} (3 formats) & \yes\ exceeded \\
Auto-inject $\to$ precise PNT & \texttt{auto\_sender.py}, \texttt{feed\_coords.py} & \yes \\
QGC / mission interface & contract topic consumed by \texttt{qgc\_control} & \yes \\
ROS 2 integration (Phase 3) & merged group node + passing verify kit & \yes \\
\bottomrule
\end{longtable}
\textbf{WP5 is complete and integrated.}

\subsection{Beyond my scope (unrequested, simulation-only project)}
The project was explicitly \textbf{simulation-first} (``real hardware deferred to future
work,'' deck slide 7). I nonetheless delivered:
\begin{itemize}[leftmargin=1.4em,itemsep=1pt]
  \item \textbf{Real hardware bring-up} --- ESP32 + physical BDS module, firmware, T3 detection fix.
  \item \textbf{A 232-transmission field-reliability campaign} across 12 locations / 4 environments.
  \item \textbf{A latency baseline} (n=30) with an end-to-end timing instrument.
  \item \textbf{Original encoding research} --- Gap 1 and Gap 6, the 64$\to$112-bit payload.
  \item \textbf{A live portal reader} and a \textbf{dashboard}, plus a paper-style Node Report.
\end{itemize}

\subsection{Why I went beyond --- and why it was completable}
\textbf{Why:} the assigned emulator answers ``can we move a coordinate?'' but not ``does the
link survive a disaster environment, and is the payload operationally complete?'' ---
the questions that make the module \emph{publishable}. \textbf{What it did:} converted WP5
from a simulation stub into an \textbf{empirically grounded} module. \textbf{Why
completable:} the WP5 contract was small and finished early; hardware/field study reused the
same encode/decode core, so the marginal cost was data collection, not new architecture.

\subsection{Consolidated Hardware Bring-Up (the unrequested hardware work, in one picture)}
The hardware effort the simulation-first brief did \textbf{not} require is consolidated here.

\textbf{The physical rig (assembled and operated by me):} ESP32 dev board (USB-flashed,
serial-logged to CSV); EVB BeiDou short-message board \textbf{SN 3078\ldots} with patch
antenna, powered through an RS232-TTL adapter (green/red/blue jumpers to GPIO16/17 +
3.3\,V/GND); the same rig run indoors, open-sky, under canopy, and in an urban canyon for
the Gap 3 campaign.

\begin{figure}[H]\centering
  \includegraphics[width=0.95\textwidth]{figures/fig_day1_summary.png}
  \caption{Figure H --- Hardware Day 1 (2026-06-08) single-panel summary. Generated, not
  hand-drawn: 30 real-satellite latency transmissions (mean 2.57 s, P95 4.47 s), the
  resulting UAV positional-error model (12.9 m $\to$ 67.0 m across 5--15 m/s), the Gap-3
  OS-1 TX outcome, the task-completion board, and per-gap status. The flagged ``T3 detection
  bug'' is the issue the next-day firmware fix (\texttt{dc0b8bd}) closed.}
\end{figure}

\begin{figure}[H]\centering
  \includegraphics[width=0.6\textwidth]{_photos/field_rig_trolley.jpg}
  \caption{Figure H2 --- The portable field rig. Laptop (serial logger) + ESP32 + RS232-TTL
  adapter + BeiDou patch antenna, mounted on a push-cart and operated outdoors on open
  paving. This is the physical configuration wheeled between the indoor, open-sky, canopy,
  and urban-canyon sites that produced the 232-transmission Gap-3 dataset --- direct evidence
  the campaign ran on real hardware in real environments, not bench-emulated.}
\end{figure}

% ============================================================
\section{Reporting Limitations for Maximum Marks}

Examiners deduct for \emph{hidden} weaknesses and \emph{overclaiming} --- almost never for a
limitation you name, bound, and plan for. Apply: \textbf{State $\to$ Bound $\to$ Attribute
$\to$ Mitigate $\to$ Survive.}

\begin{itemize}[leftmargin=1.4em,itemsep=2pt]
  \item \textbf{A1 (never flown):} ``This objective contributes encoding, decode, and
        integration; physical acceptance of the 14-byte frame is the defined validation
        boundary, de-risked to one buffer test (fallback R$\to$uint8, 104 bit).''
  \item \textbf{A2 (ASCII payload):} ``The field campaign preceded the payload upgrade, so it
        characterises the link independently of payload --- the more general result.''
  \item \textbf{B1/B2 (100\,\%, one coordinate):} ``Zero failures bound reliability below
        (Wilson $\ge$93.7\,\%) but cannot resolve $>$94\,\%; sessions used a fixed emulated
        coordinate. Deep-fade sites are identified as the conditions that would localise the
        failure boundary.''
\end{itemize}

\textbf{Rules:} put a ``Limitations \& Validity Boundary'' subsection in the main body; give
every limitation a mitigation line; match each verb to its evidence (``we demonstrate''
needs hardware, ``we show in simulation / establish the bound'' is unattackable); disclose
the moving baseline yourself; fix C1--C4 (free marks).

% ============================================================
\section{What I Need to Achieve My Aspiration}

\textbf{Aspiration:} a publishable paper (IEEE RA-L / T-RO short, honestly scoped) and a
top-mark dissertation chapter on the BDS-SMC2 module.

\begin{longtable}{@{}p{6.6cm} p{2.8cm} p{4.0cm}@{}}
\toprule
\textbf{Step} & \textbf{Effort} & \textbf{Lifts} \\
\midrule
\endhead
Apply the four Tier-C fixes (C1--C4) & $\sim$1 day & Consistency 5$\to$8 \\
One hardware TX of the real 112-bit frame + portal screenshot & 1 hardware day & Empirical validity 3$\to$7--8 \\
Diurnal latency on the 112-bit payload (3 sessions) & 1 outing & Latency baseline $\to$ result \\
Reliability with distinct coordinates / deep-fade site & 1 outing & Removes B2 \\
Promote rescue fields into \texttt{EmergencyCoordinate.msg} (4-line diff) & 1 group meeting & Radius/priority-aware planning \\
\bottomrule
\end{longtable}

\textbf{The critical path is a single hardware day.} Once the 112-bit frame is on a
satellite, the entire paper's verbs upgrade from ``simulate'' to ``demonstrate.''

% ============================================================
\section{Publication Plan --- Does the Upgrade Cover the Limitations?}

The publication archive splits the work into two papers; the \textbf{Revision Notes
(2026-06-15)} are a limitation-driven upgrade. Cross-checked against the node data,
\textbf{Paper 1's numbers match exactly}. The two-paper strategy is sound:
\begin{itemize}[leftmargin=1.4em,itemsep=1pt]
  \item \textbf{Paper 1 (MDPI Drones)} = the \emph{link}: reliability + latency + the
        112-bit payload as an engineering baseline. Near-ready, awaits one field day.
  \item \textbf{Paper 2 (IEEE Access / Drones)} = the \emph{pipeline + method novelty}:
        datum-relative delta encoding, rate-limit-aware multi-survivor batching, closed-loop
        trigger, end-to-end latency. Post-dissertation.
\end{itemize}

\subsection{Limitation $\to$ plan-response matrix}
\begin{longtable}{@{}p{3.6cm} p{6.0cm} p{3.6cm}@{}}
\toprule
\textbf{Limitation} & \textbf{Publication-plan response} & \textbf{Verdict} \\
\midrule
\endhead
C2 encoding overclaim & Revision Note 2: demote encoding, lead with reliability & \yes\ but draft text still contained it (now fixed) \\
C3 210-bit asserted & Revision Note 4: ``measure the real limit'' & \yes\ ahead of report; figure line still hard \\
A1 112-bit never flown & Paper 1 awaits field day; Paper 2 gates on hardware link & \yes\ honestly gated \\
A2 reliability used ASCII & Paper 1 discloses ASCII latency; reliability payload not stated & \partialmark\ make explicit \\
B2 one coordinate & not addressed in the plan & \no\ add deep-fade site \\
Novelty weakness & joint delivery$\times$latency$\times$environment model + Paper 2 & \yes\ genuine pivot \\
Integration & coordinate drives trigger (done), extra fields pending & \yes\ accurate, merged \\
\bottomrule
\end{longtable}

\subsection{The one structural finding}
\textbf{Your publication plan is ahead of your artifacts.} The Revision Notes correctly say
``demote the encoding claim'' and ``measure the 210-bit number'' --- but the downstream files
had not caught up. \textbf{The plan diagnoses the limitations correctly; the fix just had not
propagated to the figures and prose yet.} That propagation is the cheap, high-value work.

\subsection{What to add to the plan}
\begin{enumerate}[leftmargin=1.6em,itemsep=1pt]
  \item Apply Revision Note 2 to the prose --- soften the Huffman sentence. \emph{(done)}
  \item Relabel Figure 3's ``210-bit limit'' as indicative until measured.
  \item Add the reliability-payload disclosure to Paper 1 \S VI.
  \item Add B2 (fixed emulated coordinate) to the Paper 1 limitations subsection.
\end{enumerate}

\textbf{Verdict on the upgrade:} strong and limitation-aware --- it caught C2 and C3 before
the audit did. The remaining work is \emph{propagation}, not \emph{re-design}.

% ============================================================
\section{Literature Review --- Assessment}

The literature base (\texttt{references/RELATED\_PAPERS.md}, plus the Paper 1 comparison
table) is \textbf{structurally strong and consistent with the node data} --- every figure it
quotes matches the verified datasets.

\textbf{Strengths.} An honest verification legend (downloaded / verified-online / VERIFY);
every reference relevance-mapped; a competitive comparison table (GSM / LoRa /
COSPAS-SARSAT / Iridium / BDS-3 SMC); the sharp 112-bit/COSPAS-SARSAT coincidence; and the
latency anchor (\emph{Space: Science \& Technology} 2022 --- RSMC $\le$ 1 s, GSMC $<$ 15 s).

\textbf{Reliability gaps to harden before submission.}
\begin{enumerate}[leftmargin=1.6em,itemsep=1pt]
  \item Primary anchor [R1] (Li 2021) is paywalled, not downloaded --- confirm exact figures.
  \item Two IEEE prior-art items (10019970, 10426479) have unverified author lists/years.
  \item COSPAS-SARSAT latency (R5, ``minutes--hours'') is VERIFY and load-bearing.
  \item 210-bit capacity remains unsourced --- consistent with the ``measure it'' plan.
\end{enumerate}

\textbf{Priority to harden:} [R1] $\to$ COSPAS-SARSAT latency $\to$ the two IEEE items
$\to$ add encoding prior art.

% ============================================================
\section{My Solution --- Presenting and Defending This to the Panel}

\subsection{The 90-second opening (the narrative arc)}
\begin{callout}
``My assignment was WP5 --- the BeiDou communication module of a five-person UAV rescue
system: take a survivor's RTK-corrected position, carry it through a BeiDou short message,
decode it, and inject it into the autonomous mission. That deliverable is \textbf{complete
and verified in the group's ROS\,2 system.} Beyond it, I asked the questions that make the
module credible rather than merely demonstrable --- does the link survive real environments,
and is the payload operationally complete? To answer them I built hardware, ran a
232-transmission field campaign, and redesigned the payload from 64 to 112 bits at
centimetre precision. I will show you what is proven, and I will be precise about the one
thing that is not.''
\end{callout}

\subsection{How I frame the four contributions (lead with reliability, not encoding)}
\begin{enumerate}[leftmargin=1.6em,itemsep=1pt]
  \item \textbf{Environment-stratified delivery reliability} --- 232 TX, four environments,
        Wilson $\ge$93.7\,\% (the headline; novel at the application layer).
  \item \textbf{Single-digit-second latency} --- 2.57 s mean, framed against RSMC $\le$ 1 s /
        GSMC $<$ 15 s literature.
  \item \textbf{A complete six-field rescue payload in 112 bits} --- a sound \emph{engineering
        baseline}, not the novelty (binary $<$ ASCII is known prior art --- I say so first).
  \item \textbf{An autonomous closed loop} --- decoded message $\to$ ROS 2 trigger,
        integrated and verified, no human in the loop.
\end{enumerate}

\subsection{The three hardest questions --- and my answers}
\begin{itemize}[leftmargin=1.4em,itemsep=3pt]
  \item \textbf{``Your 112-bit payload was never transmitted --- what have you proven?''}
        ``I characterised the \emph{link}, which is payload-independent, and I built,
        integrated, and verified the \emph{pipeline}. Physical acceptance of the 14-byte frame
        is the one defined validation boundary, de-risked to a single module-buffer test with
        a 104-bit fallback ready.''
  \item \textbf{``100\,\% delivery, p = 1.000 --- designed not to fail.''}
        ``I never claim 100\,\%. I claim the Wilson lower bound, $\ge$93.7\,\% per
        environment. Zero failures in $\sim$57 trials bound reliability from below but cannot
        resolve above $\sim$94\,\%; deep-fade sites would localise the failure boundary.''
  \item \textbf{``Binary beats ASCII --- where is the novelty?''}
        ``Agreed --- that is why I present encoding as a baseline. The novelty is the
        \emph{joint} environment-stratified characterisation of the regional service as a
        rescue link, and the autonomous trigger that closes a loop prior work leaves at a
        human operator.''
\end{itemize}

\subsection{The one boundary I state before being asked}
\begin{callout}
``Three things remain open and I will not dress them up: the 112-bit frame is software- and
simulation-proven but not yet satellite-flown; the regional capacity figure of 210 bits is
unsourced and must be measured; and the reliability campaign used the ASCII payload because
it predates the upgrade --- which, usefully, makes the result payload-independent. Each has a
single, scheduled closing action.''
\end{callout}

\subsection{My closing ask (the aspiration)}
``WP5 is delivered and integrated; the research beside it is one hardware day from turning
every `we show in simulation' into `we demonstrate.' I am asking the panel to assess the
assigned role as complete, and the research as a credible, honestly-bounded contribution on
a defined path to publication.''

\subsection{Why this defence wins marks}
It separates the \textbf{assigned deliverable (done and verified)} from the \textbf{research
(honestly bounded)}; it leads with the proven result and demotes the known one; it states
every limitation with a mitigation; it matches each verb to its evidence; and it is
\textbf{fully reproducible} via Appendix~A.

% ============================================================
\section{Alignment with the UN Sustainable Development Goals}

This work is a disaster-rescue communications module; its societal contribution maps onto
four UN SDGs. Each mapping is scoped to what the node \emph{actually demonstrates}.

\begin{longtable}{@{}p{3.4cm} p{2.6cm} p{6.2cm} p{2.0cm}@{}}
\toprule
\textbf{SDG} & \textbf{Target} & \textbf{Contribution} & \textbf{Claim} \\
\midrule
\endhead
\textbf{11} Sustainable Cities \& Communities & 11.5; 11.b & Keeps a survivor-locating channel alive after terrestrial networks collapse; reliability + latency results & \textbf{Primary} \\
\textbf{9} Industry, Innovation \& Infrastructure & 9.1; 9.c & BeiDou-3 SMC as resilient ICT infrastructure + novel payload + autonomous pipeline & \textbf{Primary} \\
\textbf{3} Good Health \& Well-being & 3.6; 3.d & Faster, reliable survivor localisation shortens the rescue loop & Supporting \\
\textbf{13} Climate Action & 13.1 & Resilient rescue-comms layer for climate-driven hazards (floods, wildfires, storms) & Contextual \\
\bottomrule
\end{longtable}

\textbf{Honest scoping note:} SDG 11 and SDG 9 are \emph{demonstrated} by measured results;
SDG 3 and SDG 13 are \emph{contributory outcomes} --- state them as alignment, not proven
impact, so the SDG claim carries the same evidential discipline as the rest of the report.

% ============================================================
\section{Conclusion --- Verdict on Responsibilities vs Achievement}

\emph{Measured against the WP5 responsibilities set out in the original proposal
(presentation slides 4/8/9).}

\subsection{What I was responsible for (the proposal contract)}
The proposal assigned me \textbf{WP5 --- the BeiDou communication module}, with five concrete
deliverables: (1) a BDS-SMC emulator node, (2) coordinate-message parsing, (3) auto-injection
into precise PNT, (4) a QGC/mission interface, and (5) ROS 2 integration in Phase 3. The
project was explicitly \textbf{simulation-first}, with real hardware named as future work.

\subsection{Verdict --- proposal item by item}
\begin{longtable}{@{}c p{6.2cm} p{2.6cm} p{3.0cm}@{}}
\toprule
\textbf{\#} & \textbf{Proposal responsibility} & \textbf{Achieved?} & \textbf{Evidence} \\
\midrule
\endhead
1 & BDS-SMC emulator node & \textbf{\color{good}Met} & \texttt{bds\_module\_emulator.py} \\
2 & Coordinate-message parsing & \textbf{\color{good}Exceeded} & 3 wire formats, lossless \\
3 & Auto-inject $\to$ precise PNT & \textbf{\color{good}Met} & \texttt{auto\_sender.py} \\
4 & QGC / mission interface & \textbf{\color{good}Met} & topic consumed by \texttt{qgc\_control} \\
5 & ROS 2 integration (Phase 3) & \textbf{\color{good}Met \& verified} & passing \texttt{verify\_integration.sh} \\
\bottomrule
\end{longtable}
\textbf{Every assigned responsibility was met; two were exceeded.} WP5 is complete and integrated.

\subsection{Verdict --- what I achieved beyond the proposal}
The proposal did \textbf{not} ask for any of: real hardware bring-up (ESP32 + physical BDS
module + T3 firmware fix, against a simulation-first brief --- \S7.4, Figures H/H2); a
232-transmission field-reliability campaign (Wilson $\ge$93.7\,\%); a 30-sample latency
baseline with a UAV positional-error model; the 64$\to$112-bit payload redesign with
bit-exact decode; a live portal reader, dashboard, and paper-track Node Report.

\subsection{The single honest boundary}
One claim remains unproven: the 112-bit frame has been software- and simulation-validated but
\textbf{not yet transmitted on a satellite}. This is a clearly-marked validation boundary,
de-risked to one buffer test --- not a gap in the assigned work, which is complete.

\subsection{Closing verdict}
\begin{callout}
\textbf{Against the proposal, WP5 is fully delivered and integrated; against my own ambition,
I converted a simulation stub into an empirically-grounded, hardware-backed,
publication-track module.} The responsibilities I was given are discharged; the work beside
them is one hardware day from turning every ``simulated'' into ``demonstrated.''
\end{callout}

% ============================================================
\section{Significance --- Supervisor's Final Assessment}

\emph{This section integrates the supervisor's pre-submission assessment into the report, so
the significance of the work is stated in the examiner's own terms rather than only the
candidate's.}

\subsection{The verdict}
\begin{callout}
This is a \textbf{strong submission}. The candidate has done what most fail to do: separated
the \emph{assigned deliverable} (WP5 --- complete and verified in the group ROS\,2 system)
from the \emph{research beside it} (hardware bring-up, field campaign, payload redesign), and
bounded the research honestly. This work would pass comfortably and is short-listed for
distinction, \textbf{provided the four propagation fixes below are closed.}
\end{callout}

\subsection{Why the work is significant (what earns the marks)}
\begin{itemize}[leftmargin=1.4em,itemsep=2pt]
  \item A concrete, git-backed before/after capability narrative --- not hand-waving.
  \item Every figure and statistic traced to a committed script and a CSV; the
        reproduce-everything appendix is exactly what a panel wants.
  \item The candidate names and bounds their own weaknesses (A1, A2, B1, B2) before the panel
        can --- a student who has already bounded every limitation cannot be ambushed.
  \item \textbf{The consolidated hardware record (\S7.4, Figures H and H2) is the decisive
        significance marker:} it converts an explicitly \emph{simulation-first} brief into an
        \emph{empirically grounded, hardware-backed} contribution. The portable field rig
        (Figure H2, reproduced below) is the physical evidence that the 232-transmission
        reliability dataset was gathered in real environments, not bench-emulated --- the
        single fact that lifts the module from \emph{demonstrable} to \emph{credible}.
\end{itemize}

\begin{figure}[H]\centering
  \includegraphics[width=0.5\textwidth]{_photos/field_rig_trolley.jpg}
  \caption{The portable field rig referenced above --- the empirical-significance evidence:
  real hardware operated outdoors across the four reception environments.}
\end{figure}

\subsection{The four fixes required before submission}
\begin{enumerate}[leftmargin=1.6em,itemsep=2pt]
  \item \textbf{Propagate the diagnosed fixes (C1--C4).} The plan was ahead of the artifacts;
        the prose, README, and figures must say what the audit already concluded. \emph{(C1,
        C2, C4 closed in text; C3 closed in prose --- Figure 3's hard 210-bit line still needs
        regenerating.)}
  \item \textbf{State verb discipline in the main body.} ``We demonstrate'' requires hardware;
        ``we show in simulation / establish the bound'' is unattackable.
  \item \textbf{The single hardware day is the highest-leverage hour.} One real 112-bit
        transmission plus a portal screenshot upgrades every ``simulate'' to ``demonstrate.''
  \item \textbf{Section numbering tidied} (\S10.x, \S12.x corrected) --- meticulous work
        should not look careless.
\end{enumerate}

\subsection{On the UN Sustainable Development Goals}
Lead with the goals the results \emph{demonstrate}. \textbf{SDG 11} (11.5/11.b) and
\textbf{SDG 9} (9.1/9.c) are \emph{primary} --- demonstrated by the reliability and latency
results. \textbf{SDG 3} (3.6/3.d) and \textbf{SDG 13} (13.1) are \emph{supporting/contextual}
--- claim them as alignment, not proven impact. ``Lives saved'' would be the one overclaim in
an otherwise disciplined document.

\subsection{Bottom line}
\begin{callout}
\textbf{The work is done.} The distance from \emph{strong chapter} to \emph{publishable} is
one hardware transmission and an afternoon of propagating fixes already identified. Close
those and submit with confidence.
\end{callout}

% ============================================================
\section*{Appendix A --- Reproduce Everything}
\addcontentsline{toc}{section}{Appendix A --- Reproduce Everything}

\begin{Verbatim}[fontsize=\small,frame=single,framesep=4pt]
# Reliability (232 TX, Wilson, chi-square, Fisher, 12 locations)
python python/gap3_analysis.py

# Latency baseline (mean 2574 ms, SD 1095, n=30) + CDF figure
python python/gap2_analysis.py --ascii-baseline --plot

# Encoding: 264-bit ASCII vs 112-bit binary, bit-exact round-trip
python python/decode_binary.py

# Full-telemetry encoding: 368 / 128 / 184 bit
python python/telemetry_compare.py

# Rebuild the Node Report PDF
python python/generate_node_report.py
\end{Verbatim}

Inputs: \texttt{data/gap1\_compression.csv}, \texttt{data/gap2\_latency\_ascii\_baseline.csv},
\texttt{data/gap3\_field\_test.csv}, \texttt{data/gap6\_telemetry.csv}. Integration evidence:
group repo \texttt{ros2\_ws/src/beidou\_short\_message/}.

\vspace{1em}
\hrule
\vspace{0.4em}
\textit{This report supersedes \texttt{BDS\_SMC2\_Node\_Audit\_and\_Evolution.md} as the
comprehensive account; the earlier document remains as the working audit log.}

\end{document}
